% !TeX program = pdflatex
% Tom I: rachunek prawdopodobienstwa. Rozklady, funkcje charakterystyczne i twierdzenia graniczne.
% Plik samodzielny; kolejne czesci mozna dopisywac przed \end{document}.
\documentclass[11pt,a4paper,oneside]{report}
\usepackage[T1]{fontenc}
\usepackage[utf8]{inputenc}
\usepackage[polish]{babel}
\usepackage{lmodern}
\usepackage{amsmath,amssymb,amsthm,mathtools}
\usepackage{microtype}
\usepackage{geometry}
\usepackage{enumitem}
\usepackage{xcolor}
\usepackage{tikz}
\usetikzlibrary{arrows.meta,positioning,calc,shapes.geometric}
\usepackage[unicode]{hyperref}
% Odsylacze miedzytomowe; zbudowac wszystkie .aux i przechowywac PDF-y obok siebie.
\usepackage{xr-hyper}
\externaldocument[II-]{rachunek_prawdopodobienstwa_tom_II}[rachunek_prawdopodobienstwa_tom_II.pdf]
\externaldocument[III-]{rachunek_prawdopodobienstwa_tom_III}[rachunek_prawdopodobienstwa_tom_III.pdf]
\externaldocument[IV-]{rachunek_prawdopodobienstwa_tom_IV}[rachunek_prawdopodobienstwa_tom_IV.pdf]

\geometry{margin=2.65cm}
\hypersetup{colorlinks=true,linkcolor=blue!45!black,urlcolor=blue!45!black}
\setlength{\parindent}{0pt}
\setlength{\parskip}{4pt}
\setlength{\emergencystretch}{2em}
\setlist[enumerate]{topsep=3pt,itemsep=2pt}
\setlist[itemize]{topsep=3pt,itemsep=2pt}
\numberwithin{equation}{chapter}

\theoremstyle{plain}
\newtheorem{twierdzenie}{Twierdzenie}[chapter]
\newtheorem{lemat}[twierdzenie]{Lemat}
\newtheorem{wniosek}[twierdzenie]{Wniosek}
\newtheorem{stwierdzenie}[twierdzenie]{Stwierdzenie}
\theoremstyle{definition}
\newtheorem{definicja}{Definicja}[chapter]
\newtheorem{przyklad}{Przykład}[chapter]
\theoremstyle{remark}
\newtheorem{uwaga}{Uwaga}[chapter]
\renewcommand{\proofname}{Dowód}

\newcommand{\R}{\mathbb{R}}
\newcommand{\N}{\mathbb{N}}
\newcommand{\Z}{\mathbb{Z}}
\newcommand{\Prob}{\mathbb{P}}
\newcommand{\E}{\mathbb{E}}
\newcommand{\Var}{\operatorname{Var}}
\newcommand{\Cov}{\operatorname{Cov}}
\newcommand{\Law}{\operatorname{Law}}
\newcommand{\supp}{\operatorname{supp}}
\newcommand{\ind}[1]{\mathbf{1}_{#1}}
\newcommand{\dd}{\mathrm{d}}
\newcommand{\Bin}{\operatorname{Bin}}
\newcommand{\Poi}{\operatorname{Poi}}
\newcommand{\Geom}{\operatorname{Geom}}
\newcommand{\Unif}{\operatorname{U}}
\newcommand{\Gam}{\operatorname{Gamma}}
\newcommand{\Betad}{\operatorname{Beta}}
\newcommand{\Normal}{\mathcal{N}}
\newcommand{\iid}{\mathrel{\overset{\mathrm{iid}}{\sim}}}

\begin{document}
\begin{titlepage}
\centering
\vspace*{4cm}
{\Large Podstawy teorii prawdopodobieństwa\par}
\vspace{1.0cm}
{\Huge\bfseries Tom I. Rachunek prawdopodobieństwa\par}
\vspace{0.7cm}
{\Large Od kombinatoryki do łańcuchów Markowa i martyngałów\par}
\vfill
{\large Wersja robocza do dalszej rozbudowy\par}
\vspace{1.4cm}
\end{titlepage}
\tableofcontents

\input{../rozdzialy/rachunek_prawdopodobienstwa_tom_I/00-wprowadzenie.tex}
\input{../rozdzialy/rachunek_prawdopodobienstwa_tom_I/01-kombinatoryka-i-modele-skonczone.tex}
\input{../rozdzialy/rachunek_prawdopodobienstwa_tom_I/02-przestrzen-probabilistyczna-i-miara-prawdopodobienstwa.tex}
\input{../rozdzialy/rachunek_prawdopodobienstwa_tom_I/03-prawdopodobienstwo-klasyczne-warunkowanie-i-niezaleznosc.tex}
\input{../rozdzialy/rachunek_prawdopodobienstwa_tom_I/04-zmienne-losowe-i-rozklady-prawdopodobienstwa.tex}
\input{../rozdzialy/rachunek_prawdopodobienstwa_tom_I/05-wartosc-oczekiwana-momenty-i-podstawowe-nierownosci.tex}
\input{../rozdzialy/rachunek_prawdopodobienstwa_tom_I/06-podstawowe-rodziny-rozkladow-prawdopodobienstwa.tex}
\input{../rozdzialy/rachunek_prawdopodobienstwa_tom_I/07-zbieznosc-zmiennych-losowych-i-rozkladow.tex}
\input{../rozdzialy/rachunek_prawdopodobienstwa_tom_I/08-prawa-wielkich-liczb.tex}
\input{../rozdzialy/rachunek_prawdopodobienstwa_tom_I/09-centralne-twierdzenie-graniczne.tex}
\input{../rozdzialy/rachunek_prawdopodobienstwa_tom_I/10-warunkowa-wartosc-oczekiwana-i-projekcje-w-l2.tex}
\input{../rozdzialy/rachunek_prawdopodobienstwa_tom_I/11-dalsze-twierdzenia-graniczne-i-koncentracja-prawdopodobienstwa.tex}
\input{../rozdzialy/rachunek_prawdopodobienstwa_tom_I/12-lancuchy-markowa-w-czasie-dyskretnym.tex}
\input{../rozdzialy/rachunek_prawdopodobienstwa_tom_I/13-filtracje-adaptacja-i-czasy-zatrzymania.tex}
\input{../rozdzialy/rachunek_prawdopodobienstwa_tom_I/14-martyngaly-rachunek-strategie-i-nierownosci.tex}
\input{../rozdzialy/rachunek_prawdopodobienstwa_tom_I/15-podmartyngaly-nadmartyngaly-i-twierdzenia-o-zbieznosci.tex}
\end{document}
